\section{Related Work}
\label{sec:related}

\subsection{Vision-language GUI agents}

A series of recent systems has shown that multimodal foundation models can be turned into capable GUI agents. SeeClick~\citep{cheng2024seeclick} introduced grounding-aware pretraining for screen elements, enabling agents to localize a referent on a phone screen with high precision. CogAgent~\citep{hong2024cogagent} scaled a GUI-specialized vision-language model to high-resolution support, recognizing tiny page elements that earlier models missed. AppAgent~\citep{yang2023appagent} demonstrated an exploration-then-deployment workflow in which the agent records UI-element documentation during an offline phase and retrieves it at run time. The agent loop we use in our experiments is M3A, introduced together with AndroidWorld, which interleaves an action-selection step and a summary step so that the agent keeps a textual memory across long horizons. These works establish that vision-language models can reach competitive success rates on standard benchmarks, but they do not address the deployment-time problem of adapting to an unseen app's idiosyncrasies on a small compute budget.

\subsection{Benchmarks for mobile GUI agents}

AndroidWorld~\citep{rawles2024androidworld} is the canonical online benchmark we build on, providing 116 programmatic tasks across 20 apps with reward signals derived from the device state at episode termination. AndroidControl~\citep{li2024androidcontrol} contributes 15K demonstrations across 833 apps to study the effects of data scale on UI control. AndroidLab~\citep{xu2024androidlab} adds a reproducible benchmark and the BMOCA-style additions we use to extend our source pool. Android in the Wild~\citep{rawles2023aitw} provides a large-scale dataset of device control demonstrations. More recent benchmarks push complexity further: MobileWorld~\citep{mobileworld2025} introduces 201 tasks across 20 apps with longer horizons (27.8 vs.\ 14.3 average steps over AndroidWorld), a much higher fraction of multi-application tasks (62.2\% vs.\ 9.5\%), and additional categories covering agent--user interaction and MCP tool use. The shared limitation across all of these benchmarks is the default evaluation protocol. Training and test tasks usually share the same template family, and held-out conditions, when present, are typically defined at the parameter or seed level rather than at the template or app level. Our evaluation protocol takes the next step in this direction by partitioning at three independent levels (pool, category, template) simultaneously.

\subsection{Structured state representations for mobile agents}

A separate line of work uses finite-state machines to represent the screen-to-screen navigation structure of mobile apps. Agent-SAMA~\citep{guo2025agentsama} models app execution as an FSM whose states are UI screens and whose transitions are user actions, and builds a multi-agent system around this FSM for planning, execution verification, and error recovery, reporting 84.0\% on Mobile-Eval-E, 80.0\% on SPA-Bench, and 63.7\% on AndroidWorld. There, the FSM acts as a lightweight, model-agnostic memory layer at execution time: it is constructed online to track where the agent is in the app and to detect deviations from expected post-conditions. Our use of an FSM is different. The two-layer FSM in EvoFSM-RL is the agent's \emph{in-context knowledge}, not its execution-time state tracker; its lower layer corresponds to the app-specific navigation structure that Agent-SAMA's FSM also captures, but its upper layer encodes category-generic workflow archetypes, failure modes, and verification checklists that have no analog in Agent-SAMA's formulation. The upper layer is what we mutate online during adaptation and what transfers across apps within a category. The two FSMs are therefore complementary rather than competing: Agent-SAMA's execution-time FSM could in principle plug in as a lower-layer state tracker within our framework, with the upper-layer abstract library $L_C$ providing the category-level transfer signal it currently lacks.

\subsection{Training paradigms for GUI agents}

Two paradigms dominate the recent literature. The first relies on large-scale synthesized data and supervised fine-tuning. OpenMobile~\citep{openmobile2026} synthesizes 2.8K task instructions and 34K action steps across 20 Android apps via an exploration-then-policy-switching pipeline, then trains by SFT to reach 64.7\% on AndroidWorld with a Qwen3-VL backbone. Critically, the authors report in their appendix that adding either step-level RL (UI-R1 methodology) or trajectory-level RL (OS-Themis framework) on top of the synthesized SFT data does not yield consistent improvements over SFT alone, and they attribute this to the rich error-recovery signal already captured during synthesis. This finding holds in a large-data, in-distribution pre-deployment regime. We target the complementary regime where large-scale synthesis is infeasible at deployment time and the model must adapt online on a small sample.

The second paradigm uses online reinforcement learning to fine-tune the policy on the benchmark task distribution. MobileRL~\citep{mobilerl2025} pushes online agentic RL on mobile GUI agents to state-of-the-art success rates (80.2\% on AndroidWorld), using a difficulty-adaptive GRPO variant with positive replay, failure curriculum filtering, and shortest-path reward shaping. Their training-versus-evaluation setup places 2,000 parameter-randomized training tasks against 116 canonical test tasks drawn from the same templates, isolating the effect of distributional augmentation rather than cross-template or cross-app generalization. Our evaluation explicitly disjoints templates between $T_{\text{adapt}}$ and $T_{\text{eval}}$ and disjoints apps between source and target pools, so absolute numbers are not directly comparable; the work is methodologically complementary in that MobileRL studies algorithmic modifications to GRPO under in-distribution evaluation while we study orthogonal context-evolution as an adaptation axis under strict cross-distribution evaluation.

\subsection{Test-time adaptation and meta-learning}

Test-time adaptation has a long line of work in vision~\citep{sun2020ttt,wang2021tent} arguing that a model should update itself online on the distribution it is asked to predict, often using self-supervised objectives or entropy minimization. These methods assume an unsupervised stream of test data and target image classification; they do not directly apply to multi-step agentic settings where reward is sparse and trajectories are long, but they motivate the deployment-time adaptation framing we adopt. Closer in spirit is the meta-learning literature on adapting from a few labeled examples per task~\citep{finn2017maml}, which prescribes a strict support/query partition we adopt as our $T_{\text{adapt}}$/$T_{\text{eval}}$ split. The generalization-in-RL literature~\citep{cobbe2020procgen} similarly argues for explicit train/test level separation rather than parameter randomization, an argument we extend to the cross-app and cross-category levels.

\subsection{Joint context and weight optimization}

The most direct precursor of our method is E-SPL~\citep{zhang2026espl}, which jointly evolves a free-text system prompt by structured mutation under a frozen reference LLM while simultaneously updating model weights via group-relative policy optimization. E-SPL operates in the single-turn reasoning setting (math, agentic web search) where rollouts are cheap and rewards are densely defined. We port the joint paradigm to the multi-step, visual, expensive-rollout GUI agent setting, with three concrete changes. The free-text prompt becomes a structured two-layer FSM with a linter that enforces transferability of the upper layer. The within-prompt advantage baseline becomes a within-(FSM, task) baseline that accommodates the fact that we run far fewer rollouts per cell. And the same joint loop runs at both source-pool pretraining and target-app deployment time, with a category-level abstract library $L_C$ mediating between them.

\subsection{Foundation models and parameter-efficient fine-tuning}

We use Qwen3-VL~\citep{qwen3vl2025} as our vision-language backbone, fine-tuning only via LoRA adapters~\citep{hu2021lora} on the attention projections at rank 16. LoRA has become standard for sample-efficient adaptation of large pretrained models and fits our deployment-time compute budget of roughly 80 trajectories per target app. Our weight update is group-relative policy optimization~\citep{shao2024deepseekmath} from DeepSeekMath, modified only in the group key (within-(FSM, task) rather than within-prompt) to match our setting. We use TrueSkill~\citep{herbrich2007trueskill} as the rating system that drives population selection during in-context evolution, with optimistic softmax sampling that keeps newly introduced variants in contention before they have accumulated evidence.

\subsection{Summary of differentiation}

Three axes summarize where this paper sits relative to recent work. First, on \textbf{evaluation strictness}: we evaluate cross-app, cross-category, and cross-template generalization simultaneously, while dominant prior protocols evaluate within-template parameter robustness. Second, on \textbf{mechanism}: we treat structured in-context knowledge (FSM) and policy weights (LoRA) as two co-adapted axes rather than treating context as a fixed prompt or weights as the only learnable component. Third, on \textbf{deployment-time framing}: we adapt under a small sample budget on a held-out target app, separating this from the large-data fine-tuning regimes that drive most current AndroidWorld numbers.
